\documentclass{alex_summary}

\begin{document}
\section*{Reading summary}
In \blockcquote{giovannetti_quantum-enhanced_2004}{Quantum-Enhanced Measurements: Beating the Standard Quantum Limit} the authors outline strategies to \textbf{overcome precision limitations} when measuring using \blockquote{classical} techniques. The fundamental limit to measuring accurately, the \textbf{Heisenberg limit}, is well below the \textbf{shot noise or standard quantum limit} that conventional measurements struggle with. By leveraging \textbf{quantum correlations} in measurement procedures the measurement uncertainty can be lowered.

The conventional way of reducing statistical measurement uncertainties relies on \textbf{repeated measurement}. The relative uncertainty of an observable scales with $\mathbf{1/\sqrt{N}}$, with \( N \) the number of performed measurements. This scaling behaviour is not ideal, as one quickly reaches a point of diminishing return and thus reducing the relative error is only practical up to a few orders of magnitude. Using \blockquote{quantum tricks} as the authors nicely phrase it, this scaling can be \textbf{improved} to a \( \mathbf{1/N} \) dependence, giving a more direct benefit from repeated measurements. 

An interesting discussion arose during the end of the paper, in which the lower limit to \textbf{precise measurements of spacetime geometry} is analyzed. Resolving \textbf{small scales} involves accurate knowledge of time and space. Both can be improved by \textbf{increasing the energy} of the probe, increasing the \blockquote{clock tick length} and deBroglie wavelength respectively. This can be done until one reaches energies close to the \textbf{Planck scale}, at which point our current knowledge breaks down and a \textbf{quantum gravitational treatment} is required. Mapping out \textbf{large-scale} spacetime geometry can be done by spreading out a \textbf{swarm of ticking clocks} that exchange and compare signals between each other. By associating a tick event with energy one can derive a \textbf{\blockquote{clock-tick} density} at which black hole formation occurs, rendering the entire experiment infeasible. This limit implies a \textbf{trade-off} between \textbf{temporal} (sparse, fast-ticking clocks) and \textbf{spatial} (dense, slow-ticking clocks) \textbf{resolution}, which is consistent with other means of bounding the accuracy of mapping out spacetime like the holographic bound, Bekenstein bound or the covariant entropy bound. \\

\begin{question1}
	What is a biphoton? Is it just two photons described by a single wavefunction or do they have to fulfill certain requirements?
\end{question1}

\begin{question3}
	With what fidelity can photonic NOON states be created? Is it large enough to see improvements in measurement precision or does the added infidelity of initializing more complicated states compensate its benefit?
\end{question3}

\printbibliography
\end{document}
